Beide Seiten einer Gleichung haben dieselbe Einheit, ebenso alle Summanden.
\begin{abcliste}
\abc $p = a\,t$ hat die Einheit $\si{\meter\per\second\squared}\cdot\si{\second}$, also $[p] = \result{\si{\meter\per\second}}$.
\abc $a\,p$ ist eine Geschwindigkeit: $[p] = \dfrac{\si{\meter\per\second}}{\si{\meter\per\second\squared}} = \result{\si{\second}}$. Der Summand $q$ ist ebenfalls eine Geschwindigkeit: $[q] = \result{\si{\meter\per\second}}$.
\abc $p = \dfrac{a}{s^2}$ hat die Einheit $\dfrac{\si{\meter\per\second\squared}}{\si{\meter\squared}}$, also $[p] = \result{\si{\per\meter\per\second\squared}}$.
\abc $p = \dfrac{v^2}{a\,s}$ hat die Einheit $\dfrac{\si{\meter\squared\per\second\squared}}{\si{\meter\per\second\squared}\cdot\si{\meter}} = 1$: $p$ ist $\result{\text{dimensionslos}}$.
\abc $p\,v$ ist eine Beschleunigung: $[p] = \dfrac{\si{\meter\per\second\squared}}{\si{\meter\per\second}} = \result{\si{\per\second}}$. Ebenso $\dfrac{s}{q}$: $[q] = \dfrac{\si{\meter}}{\si{\meter\per\second\squared}} = \result{\si{\second\squared}}$.
\end{abcliste}
